% Motor: XeLaTeX, fijado en .latexmkrc y en la receta del editor.
\documentclass[aspectratio=169,10pt,c]{beamer}

% !TEX root = slides.tex
% Preambulo comun: estilo UTP, colores, fuentes, plantillas de beamer y
% comandos de apoyo didactico. Se carga desde slides.tex tras \documentclass.

% La opción "c" de beamer reparte el espacio libre 1:1.5 (arriba:abajo), por lo
% que el contenido queda por encima del centro. Igualamos ambos rellenos.
\makeatletter
\define@key{beamerframe}{c}[true]{%
  \beamer@frametopskip=0pt plus 1fill\relax%
  \beamer@framebottomskip=0pt plus 1fill\relax%
  \beamer@frametopskipautobreak=0pt plus .5\paperheight\relax%
  \beamer@framebottomskipautobreak=0pt plus .5\paperheight\relax%
}
\makeatother

\usepackage{graphicx}
\usepackage{tikz}
\usepackage{xcolor}
\usepackage{iftex}

\ifPDFTeX
  \usepackage[utf8]{inputenc}
  \usepackage[T1]{fontenc}
  \usepackage{lmodern}          % tipografia escalable para el respaldo con pdflatex
  \renewcommand{\familydefault}{\sfdefault}
\else
  \usepackage{fontspec}
  \IfFontExistsTF{Arial}
    {\setsansfont{Arial}[BoldFont={Arial Bold}]}
    {\setsansfont{Helvetica}[BoldFont={Helvetica Bold}]}
\fi

\usepackage[spanish,es-nodecimaldot]{babel}
\usepackage{microtype}
\usepackage{booktabs}
\usepackage{amssymb}
\usepackage{array}
\usetikzlibrary{arrows.meta,positioning,fit,calc,shapes.geometric}

\graphicspath{{assets/}}

% ---------- Identidad visual UTP ----------
\definecolor{utpRojo}{HTML}{C42B35}
\definecolor{utpNegro}{HTML}{111111}
\definecolor{utpGris}{HTML}{5F5F5F}
\definecolor{utpGrisMedio}{HTML}{9A9A9A}
\definecolor{utpGrisClaro}{HTML}{F1F1F1}
\definecolor{utpNaranja}{HTML}{C1670B}

\usetheme{default}
\setbeamercolor{normal text}{fg=utpNegro,bg=white}
\setbeamercolor{structure}{fg=utpRojo}
\setbeamercolor{frametitle}{fg=utpNegro,bg=white}
\setbeamercolor{title}{fg=utpNegro}
\setbeamercolor{subtitle}{fg=utpNegro}
\setbeamercolor{author}{fg=utpNegro}
\setbeamercolor{date}{fg=utpNegro}
\setbeamercolor{block title}{fg=white,bg=utpNegro}
\setbeamercolor{block body}{fg=utpNegro,bg=utpGrisClaro}
\setbeamercolor{alertblock title}{fg=white,bg=utpRojo}
\setbeamercolor{alertblock body}{fg=utpNegro,bg=utpRojo!7}
\setbeamercolor{exampleblock title}{fg=white,bg=utpGris}
\setbeamercolor{exampleblock body}{fg=utpNegro,bg=utpGris!8}

\setbeamerfont{title}{series=\bfseries,size=\fontsize{19.5}{22.5}\selectfont}
\setbeamerfont{subtitle}{size=\fontsize{11.5}{13.5}\selectfont}
\setbeamerfont{date}{series=\bfseries,size=\fontsize{11.5}{13.5}\selectfont}
\setbeamerfont{author}{size=\fontsize{10.5}{12.5}\selectfont}
\setbeamerfont{frametitle}{series=\bfseries,size=\large}
\setbeamerfont{block title}{series=\bfseries,size=\small}
\setbeamerfont{block body}{size=\normalsize}

\setbeamertemplate{navigation symbols}{}
\setbeamertemplate{headline}{}
\setbeamertemplate{blocks}[default]
\setbeamertemplate{itemize item}{\color{utpRojo}\small\raise.15ex\hbox{$\blacktriangleright$}}
\setbeamertemplate{itemize subitem}{\color{utpGris}\scriptsize$\bullet$}
\setbeamertemplate{enumerate item}{\color{utpRojo}\insertenumlabel.}
\setbeamertemplate{background canvas}{\color{white}\rule{\paperwidth}{\paperheight}}

% Título de diapositiva: negrita negra a la izquierda, logo UTP a la derecha
\setbeamertemplate{frametitle}{%
  \begin{tikzpicture}[remember picture,overlay]
    \node[anchor=north east,inner sep=0pt]
      at ([xshift=-7mm,yshift=-5.5mm]current page.north east)
      {\includegraphics[width=17mm]{utp-logo.png}};
  \end{tikzpicture}%
  \vspace{1.2mm}%
  \begin{minipage}[t]{.80\textwidth}\raggedright\insertframetitle\end{minipage}\par
  \vspace{1.1mm}{\color{utpRojo}\rule{\textwidth}{1.1pt}}
}

% Solo numeración discreta; la plantilla UTP no lleva pie de página
\setbeamertemplate{footline}{%
  \hfill{\color{utpGrisMedio}\scriptsize\insertframenumber}\hspace{6mm}\vspace{2.5mm}%
}

% Portada UTP
\setbeamertemplate{title page}{%
  \begin{tikzpicture}[remember picture,overlay]
    \node[anchor=center,inner sep=0pt] at (current page.center) {%
      \includegraphics[width=\paperwidth,height=\paperheight]{utp-cover-bg.jpeg}%
    };
    \node[anchor=north,align=center,text=utpNegro,inner sep=0pt,
      text width=145mm,font=\usebeamerfont{title}]
      at ([yshift=-25.2mm]current page.north) {\inserttitle};
    \node[anchor=north,align=center,text=utpNegro,inner sep=0pt,
      text width=142mm,font=\usebeamerfont{subtitle}]
      at ([yshift=-38.2mm]current page.north) {\insertsubtitle};
    \node[anchor=north,align=center,text=utpNegro,inner sep=0pt,
      font=\usebeamerfont{date}]
      at ([yshift=-50.4mm]current page.north) {\insertdate};
    \node[anchor=north,align=center,text=utpNegro,inner sep=0pt,
      font=\usebeamerfont{author}]
      at ([yshift=-58.4mm]current page.north) {\insertauthor};
  \end{tikzpicture}%
}

% ---------- Comandos de apoyo didáctico ----------
\newcommand{\concepto}[2]{%
  \begin{block}{#1}\large #2\end{block}}
\newcommand{\pregunta}[1]{%
  \begin{alertblock}{Decide y explica}#1\end{alertblock}}
\newcommand{\analogia}[1]{%
  \begin{exampleblock}{En palabras simples}\small #1\end{exampleblock}}
\newcommand{\enlarealidad}[1]{%
  \begin{block}{Cómo se ve en la realidad}\small #1\end{block}}
\newcommand{\ojo}[1]{%
  \begin{alertblock}{Ojo con esto}\small #1\end{alertblock}}
\newcommand{\apoyo}{{\normalfont\footnotesize\color{utpGris}~· material de apoyo}}
\newcommand{\marca}[1]{\tikz[baseline=(m.base)]{%
  \node[rounded corners=1.5pt,fill=utpGrisClaro,draw=utpGrisMedio,inner xsep=3pt,inner ysep=1.5pt,
        font=\footnotesize] (m) {#1};}}

\tikzset{
  caja/.style={draw=utpNegro, very thick, rounded corners=1.5pt, fill=white,
    align=center, minimum height=0.72cm, inner sep=5pt, font=\small},
  cajaazul/.style={caja, draw=utpGris, fill=white},
  cajaverde/.style={caja, draw=utpNegro, fill=utpGrisClaro},
  cajanaranja/.style={caja, draw=utpRojo!75, fill=utpRojo!7},
  flecha/.style={-{Stealth[length=2.3mm]}, very thick, draw=utpNegro},
  flechavuelta/.style={{Stealth[length=2.3mm]}-{Stealth[length=2.3mm]}, very thick, draw=utpNegro}
}

% ---------- Notación pata de gallo (crow's foot) ----------
% La rotación indica hacia dónde está la entidad: 0 derecha, 180 izquierda,
% 90 arriba, -90 abajo. La pata se abre hacia la entidad del lado "muchos".
\tikzset{
  relacion/.style={thick, draw=utpNegro},
  pics/uno/.style={code={%
    \draw[thick,utpNegro,line cap=round] (0,-1.5mm) -- (0,1.5mm);}},
  pics/muchos/.style={code={%
    \draw[thick,utpNegro,line cap=round,line join=round]
      (2.6mm,1.5mm) -- (0,0) -- (2.6mm,-1.5mm);
    \draw[thick,utpNegro,line cap=round] (0,0) -- (2.6mm,0);}},
}
\newcommand{\simbuno}{\tikz[baseline=-0.6ex]{\draw[relacion](0,0)--(7mm,0);\pic at (4.4mm,0){uno};}}
\newcommand{\simbmuchos}{\tikz[baseline=-0.6ex]{\draw[relacion](0,0)--(7mm,0);\pic at (4.4mm,0){muchos};}}

\usepackage[normalem]{ulem}

\title{Índices, nulos, expresiones y agregaciones}
\subtitle{Buscar más rápido, tratar los datos que faltan, calcular y resumir}
\date{Semana 09 - Sesión 18}
\author{Rel Guzman}

% Espacio en blanco para completar
\newcommand{\bl}[1][1.9cm]{\underline{\hspace{#1}}}
% Nombre de tabla, siempre en naranja
\newcommand{\tbl}[1]{\textcolor{utpNaranja}{\texttt{#1}}}
% Palabra reservada de SQL dentro de un bloque de código
\newcommand{\kw}[1]{\textcolor{utpRojo}{\textbf{#1}}}
% Parte nueva de un bloque de codigo que crece paso a paso
\newcommand{\nuevo}[1]{{%
  \setlength{\fboxsep}{1.2pt}%
  \begingroup
  \renewcommand{\kw}[1]{\textcolor{utpNegro}{\textbf{##1}}}%
  \colorbox{yellow!45}{\textcolor{utpNegro}{#1}}%
  \endgroup}}
% Comentario dentro de un bloque de código SQL
\newcommand{\com}[1]{\textcolor{utpGris}{-- #1}}
% Marcas de clave primaria y clave foránea
\newcommand{\pk}{{\scriptsize\color{utpRojo}\textbf{PK}}}
\newcommand{\fk}{{\scriptsize\color{utpNaranja}\textbf{FK}}}
% Insignia de práctica
\newcommand{\ins}[1]{\tikz[baseline=-0.5ex]{\node[rounded corners=2pt,fill=utpRojo,text=white,
  font=\bfseries\scriptsize,inner xsep=5pt,inner ysep=2.5pt] {#1};}}
% Etiqueta oscura: nombre del gestor o de la forma de escribir algo
\newcommand{\etiqueta}[1]{\tikz[baseline=-0.5ex]{\node[rounded corners=2pt,fill=utpNegro,text=white,
  font=\bfseries\scriptsize,inner xsep=6pt,inner ysep=3pt] {#1};}}
% Celda de una fila que se quedó sin pareja
\newcommand{\mal}[1]{{\setlength{\fboxsep}{1.2pt}\colorbox{utpRojo!18}{#1}}}
% Valor desconocido
\newcommand{\nulo}{\textcolor{utpGrisMedio}{\textit{NULL}}}
% Bloque de código sobre fondo gris
\newsavebox{\sqlbuf}
\newenvironment{sqlcode}[1][.86\textwidth]
  {\begin{center}\begin{lrbox}{\sqlbuf}\begin{minipage}{#1}\ttfamily\footnotesize\vspace{6pt}\raggedright}
  {\vspace{6pt}\end{minipage}\end{lrbox}\colorbox{utpGrisClaro}{\usebox{\sqlbuf}}\end{center}}
% Sangría dentro de un bloque de código
\newcommand{\ind}{\hspace*{1.6em}}
% Anatomía de una instrucción: cada pieza del código con su pregunta al lado
\tikzset{
  pieza/.style={fill=utpGrisClaro,font=\ttfamily\footnotesize,anchor=west,inner sep=5pt,
    minimum height=0.62cm},
  para/.style={font=\small,anchor=west,text=utpNegro}
}
\newcommand{\piezax}{7.6}
\newcommand{\pieza}[3]{%
  \node[pieza] at (0,#1) {#2};
  \node[para] at (\piezax,#1) {\textcolor{utpRojo}{$\blacktriangleleft$}\ #3};}
% Diapositiva de sección: una palabra bien grande (\seccion[0]{...} la deja sin línea roja)
\newcommand{\seccion}[2][1]{%
\begin{frame}[plain]
  \begin{tikzpicture}[remember picture,overlay]
    \node[anchor=north east,inner sep=0pt]
      at ([xshift=-10mm,yshift=-8mm]current page.north east)
      {\includegraphics[width=43mm]{utp-logo.png}};
    \node[anchor=center,inner sep=0pt] at (current page.center)
      {\resizebox{100mm}{!}{\bfseries\color{utpNegro}#2}};
    \ifnum#1=1
    \node[anchor=north,inner sep=0pt] at ([yshift=-11mm]current page.center)
      {\color{utpRojo}\rule{100mm}{1.6pt}};
    \fi
  \end{tikzpicture}
\end{frame}}
% Diagrama de dos círculos: se pinta la parte que el JOIN conserva.
% \venn[escala]{izquierda}{centro}{derecha}, cada uno con \von o \voff
\newcommand{\von}{utpRojo!45}
\newcommand{\voff}{white}
\newcommand{\venn}[4][1]{%
  \begin{tikzpicture}[scale=#1,baseline=-0.6ex]
    \fill[#2] (0,0) circle (0.62);
    \fill[#4] (0.74,0) circle (0.62);
    \begin{scope}
      \clip (0,0) circle (0.62);
      \fill[#3] (0.74,0) circle (0.62);
    \end{scope}
    \draw[thick,utpNegro] (0,0) circle (0.62);
    \draw[thick,utpNegro] (0.74,0) circle (0.62);
  \end{tikzpicture}}
% Cabecera de columna dentro de una tabla en monoespaciado
\newcommand{\cab}[1]{\normalfont\bfseries #1}
% Resultado de una consulta: \res{columnas}{cabecera}{filas}
\newcommand{\res}[3]{%
  \centering{\scriptsize\color{utpGris}resultado}\\[3pt]
  \footnotesize\ttfamily
  \begin{tabular}{#1}
  \toprule
  #2 \\
  \midrule
  #3
  \bottomrule
  \end{tabular}}
% Tablas de la revisión: la tienda
\newcommand{\tablacliente}{%
  \centering\tbl{cliente}\\[4pt]
  \scriptsize\ttfamily
  \begin{tabular}{rl}
  \toprule
  \cab{id} & \cab{nombre} \\
  \midrule
  1 & Ana \\
  2 & Bruno \\
  3 & Carla \\
  \bottomrule
  \end{tabular}}
\newcommand{\tablapedido}{%
  \centering\tbl{pedido}\\[4pt]
  \scriptsize\ttfamily
  \begin{tabular}{rl}
  \toprule
  \cab{cliente\_id} & \cab{producto} \\
  \midrule
  1 & Laptop \\
  1 & Mouse \\
  3 & Teclado \\
  4 & Monitor \\
  \bottomrule
  \end{tabular}}
% Tabla de ejemplo de la sesión: los empleados
\newcommand{\tablaempleado}{%
  \centering\tbl{empleado}\\[4pt]
  \scriptsize\ttfamily
  \setlength{\tabcolsep}{4pt}
  \begin{tabular}{rllrrl}
  \toprule
  \cab{id} & \cab{nombre} & \cab{area} & \cab{sueldo} & \cab{anio\_ingreso} & \cab{telefono} \\
  \midrule
  1 & Ana & Ventas & 2500 & 2018 & 987111222 \\
  2 & Bruno & Ventas & 1800 & 2021 & \nulo \\
  3 & Carla & Sistemas & 4200 & 2015 & 955333444 \\
  4 & Diego & Sistemas & 3100 & 2023 & '' \\
  5 & Elena & Almacén & 1500 & 2020 & \nulo \\
  \bottomrule
  \end{tabular}}
% Tabla de práctica: los productos
\newcommand{\tablaproducto}{%
  \centering\tbl{producto}\\[4pt]
  \scriptsize\ttfamily
  \setlength{\tabcolsep}{4pt}
  \begin{tabular}{rllrr}
  \toprule
  \cab{id} & \cab{nombre} & \cab{categoria} & \cab{precio} & \cab{stock} \\
  \midrule
  1 & Laptop & Cómputo & 3200 & 5 \\
  2 & Mouse & Cómputo & 45 & \nulo \\
  3 & Silla & Muebles & 380 & 12 \\
  4 & Mesa & Muebles & 650 & 0 \\
  5 & Lámpara & Hogar & 90 & \nulo \\
  \bottomrule
  \end{tabular}}

\begin{document}

% 1 ---------------------------------------------------------------------
{\setbeamertemplate{footline}{}\maketitle}

% 2 ---------------------------------------------------------------------
\begin{frame}[c]
\begin{block}{Logro de la sesión}
Al finalizar, podrás explicar \textbf{qué es un índice}, tratar los \textbf{valores nulos}, calcular con \textbf{expresiones} y resumir datos con \textbf{agregaciones}.
\end{block}
\end{frame}
% =======================================================================
% Revisión de la sesión anterior
% ---------------------------------------------------------------------
\begin{frame}{Revisión de la sesión anterior: ¿qué hace un \texttt{JOIN}?}
\begin{columns}[c,onlytextwidth]
\column{.20\textwidth}
\tablacliente\\[3pt]
{\normalfont\scriptsize\color{utpGris}tabla \textbf{principal}}
\column{.28\textwidth}
\tablapedido\\[3pt]
{\normalfont\scriptsize\color{utpGris}la \textbf{otra} tabla}
\column{.48\textwidth}
\small Tienes una \textbf{tabla principal} y buscas \textbf{coincidencias} en otra tabla, usando una \textbf{columna en común}.
\begin{sqlcode}[\textwidth]
\scriptsize
\kw{ON} cliente.id = pedido.cliente\_id
\end{sqlcode}
\end{columns}
\vspace{0.25cm}
\begin{itemize}\setlength\itemsep{5pt}
  \item \textbf{Ana} tiene dos coincidencias y \textbf{Carla} una.
  \item \mal{Bruno} no tiene ninguna: no hizo pedidos.
  \item El \mal{Monitor} tampoco: su \texttt{cliente\_id} es 4 y ese cliente no existe.
\end{itemize}
\vspace{0.2cm}
\concepto{La pregunta de siempre}{\small Cada tipo de \texttt{JOIN} decide \textbf{qué hacer con las filas que no encontraron coincidencia}.}
\end{frame}
% ---------------------------------------------------------------------
\begin{frame}{Revisión de la sesión anterior: los cuatro \texttt{JOIN}}
\begin{center}
\renewcommand{\arraystretch}{1.7}
\small
\begin{tabular}{>{\centering\arraybackslash}m{1.5cm}>{\ttfamily\bfseries\raggedright\arraybackslash}m{2.9cm}>{\raggedright\arraybackslash}m{3.5cm}>{\raggedright\arraybackslash}m{4.9cm}}
\toprule
& \normalfont\bfseries Tipo & \bfseries Qué conserva & \bfseries En el ejemplo \\
\midrule
\venn[0.42]{\voff}{\von}{\voff} & INNER JOIN & solo las coincidencias & Ana y Carla. Bruno no aparece. \\
\venn[0.42]{\von}{\von}{\voff} & LEFT JOIN & toda la tabla izquierda & Ana, Bruno y Carla. Bruno con \nulo. \\
\venn[0.42]{\voff}{\von}{\von} & RIGHT JOIN & toda la tabla derecha & todos los pedidos, también el Monitor. \\
\venn[0.42]{\von}{\von}{\von} & FULL OUTER JOIN & todo, de ambas tablas & todos; lo que no coincide va con \nulo. \\
\bottomrule
\end{tabular}
\end{center}
\vspace{0.15cm}
\begin{columns}[T,onlytextwidth]
\column{.32\textwidth}
\centering\etiqueta{LEFT JOIN}\\[5pt]
\small «mi tabla principal \textbf{manda}»
\column{.32\textwidth}
\centering\etiqueta{INNER JOIN}\\[5pt]
\small «solo me interesan las \textbf{coincidencias}»
\column{.32\textwidth}
\centering\etiqueta{FULL JOIN}\\[5pt]
\small «quiero \textbf{todo}, coincida o no»
\end{columns}
\end{frame}
% ---------------------------------------------------------------------
\begin{frame}{Revisión de la sesión anterior: leer un \texttt{LEFT JOIN}}
\begin{columns}[c,onlytextwidth]
\column{.52\textwidth}
\begin{sqlcode}[\textwidth]
\scriptsize
\kw{SELECT} c.id, c.nombre, p.producto\\
\kw{FROM} cliente c\\
\nuevo{\kw{LEFT JOIN}} pedido p\\
\ind \kw{ON} c.id = p.cliente\_id;
\end{sqlcode}
\column{.42\textwidth}
\res{rll}{\cab{id} & \cab{nombre} & \cab{producto}}{%
  1 & Ana & Laptop \\
  1 & Ana & Mouse \\
  2 & \mal{Bruno} & \mal{\nulo} \\
  3 & Carla & Teclado \\}
\end{columns}
\vspace{0.2cm}
\analogia{«Mi tabla principal es \tbl{cliente}. Quiero conservar \textbf{todos} los clientes y buscar coincidencias en \tbl{pedido}. Si las encuentro, agrego esos datos. Si no, pongo \nulo.»}
\vspace{0.05cm}
\ojo{Una fila de la tabla principal \textbf{se repite} si tiene varias coincidencias: Ana sale \textbf{dos veces} porque tiene dos pedidos.}
\end{frame}
% =======================================================================
\seccion[0]{Índices}
% ---------------------------------------------------------------------
\begin{frame}{Sin índice: buscar fila por fila}
\begin{columns}[c,onlytextwidth]
\column{.30\textwidth}
\centering\tbl{cliente}\\[4pt]
\scriptsize\ttfamily
\begin{tabular}{rl}
\toprule
\cab{id} & \cab{nombre} \\
\midrule
1 & Rosa \\
2 & Beto \\
3 & Zoe \\
4 & Ana \\
5 & Luis \\
\bottomrule
\end{tabular}
\column{.66\textwidth}
\begin{sqlcode}[\textwidth]
\scriptsize
\kw{SELECT} * \kw{FROM} cliente\\
\kw{WHERE} nombre = 'Luis';
\end{sqlcode}
\vspace{-0.1cm}
\small
\begin{itemize}\setlength\itemsep{6pt}
  \item Los nombres están \textbf{desordenados}.
  \item Hay que preguntar \textbf{fila por fila}: ¿Rosa? no. ¿Beto? no. ¿Zoe? no\ldots
  \item Con 5 filas no se nota. Con \textbf{un millón}, sí.
\end{itemize}
\end{columns}
\vspace{0.5cm}
\analogia{Es buscar una palabra en un libro \textbf{leyéndolo completo}, página por página.}
\end{frame}
% ---------------------------------------------------------------------
\begin{frame}{Con índice: ir directo a la fila}
\begin{columns}[c,onlytextwidth]
\column{.52\textwidth}
\centering
\begin{tikzpicture}[y=0.56cm]
  \tikzset{
    celda/.style={draw=utpNegro,thick,font=\scriptsize\ttfamily,minimum height=0.48cm,
      inner sep=2.5pt,align=left},
    idx/.style={celda,fill=yellow!35,text width=1.15cm},
    dato/.style={celda,fill=utpGrisClaro,text width=1.75cm}
  }
  \node[idx] (i1) at (0,0) {Ana};
  \node[idx] (i2) at (0,-1) {Beto};
  \node[idx] (i3) at (0,-2) {Luis};
  \node[idx] (i4) at (0,-3) {Rosa};
  \node[idx] (i5) at (0,-4) {Zoe};
  \node[dato] (t1) at (4.3,0) {1 \ Rosa};
  \node[dato] (t2) at (4.3,-1) {2 \ Beto};
  \node[dato] (t3) at (4.3,-2) {3 \ Zoe};
  \node[dato] (t4) at (4.3,-3) {4 \ Ana};
  \node[dato] (t5) at (4.3,-4) {5 \ Luis};
  \draw[-{Stealth[length=1.8mm]},utpGris,thick] (i1.east)--(t4.west);
  \draw[-{Stealth[length=1.8mm]},utpGris,thick] (i2.east)--(t2.west);
  \draw[-{Stealth[length=1.8mm]},utpRojo,very thick] (i3.east)--(t5.west);
  \draw[-{Stealth[length=1.8mm]},utpGris,thick] (i4.east)--(t1.west);
  \draw[-{Stealth[length=1.8mm]},utpGris,thick] (i5.east)--(t3.west);
  \node[font=\scriptsize\bfseries,above=1mm of i1] {índice};
  \node[font=\scriptsize,above=1mm of t1] {\tbl{cliente}};
  \node[font=\tiny,text=utpGris,below=1mm of i5] {ordenado};
\end{tikzpicture}
\column{.46\textwidth}
\begin{itemize}\setlength\itemsep{9pt}
  \item Un índice es una \textbf{lista ordenada} de una columna.
  \item Cada nombre \textbf{señala} dónde está su fila.
  \item La tabla \textbf{queda igual}.
\end{itemize}
\end{columns}
\vspace{0.4cm}
\analogia{Es el \textbf{índice de un libro}: buscas la palabra en la lista ordenada y te dice en qué página está.}
\end{frame}
% ---------------------------------------------------------------------
\begin{frame}{\texttt{CREATE INDEX}: crear un índice}
\begin{center}
\begin{tikzpicture}[y=0.80cm]
  \renewcommand{\piezax}{6.0}
  \pieza{1}{\nuevo{\kw{CREATE INDEX}} idx\_cliente\_nombre}{el \textbf{nombre} que le pones al índice}
  \pieza{0}{\nuevo{\kw{ON}} cliente (nombre);}{la \textbf{tabla} y la \textbf{columna}}
\end{tikzpicture}
\end{center}
\vspace{0.3cm}
\begin{itemize}\setlength\itemsep{8pt}
  \item Se crea \textbf{una sola vez}.
  \item Tus consultas \textbf{se escriben igual}: solo salen más rápido.
  \item Sirve en las columnas por las que \textbf{buscas seguido}.
\end{itemize}
\vspace{0.3cm}
\ojo{No se le pone índice a todo: cada índice \textbf{ocupa espacio} y hace un poco más lento \textbf{guardar} datos nuevos.}
\end{frame}
% ---------------------------------------------------------------------
\begin{frame}{Demostración: buscar entre un millón de contactos}
\small La carpeta \texttt{indices/} trae \textbf{la misma agenda dos veces}. Solo cambia \textbf{una instrucción}.
\vspace{0.15cm}
\begin{columns}[T,onlytextwidth]
\column{.48\textwidth}
\centering\etiqueta{sin-indice}
\begin{sqlcode}[\textwidth]
\scriptsize
\kw{CREATE TABLE} contacto (\\
\ind id \ \ \ \ \ INTEGER \kw{PRIMARY KEY},\\
\ind nombre \ TEXT,\\
\ind correo \ TEXT\\
);
\end{sqlcode}
\column{.48\textwidth}
\centering\etiqueta{con-indice}
\begin{sqlcode}[\textwidth]
\scriptsize
\com{la misma tabla, y además:}\\[4pt]
\nuevo{\kw{CREATE INDEX} idx\_contacto\_correo}\\
\nuevo{\ind \kw{ON} contacto (correo);}
\end{sqlcode}
\end{columns}
\vspace{0.15cm}
\begin{center}
\small
\renewcommand{\arraystretch}{1.3}
\begin{tabular}{>{\raggedright\arraybackslash}p{6.0cm}>{\raggedright\arraybackslash}p{2.9cm}>{\raggedright\arraybackslash}p{2.9cm}}
\toprule
& \bfseries Sin índice & \bfseries Con índice \\
\midrule
Buscar un contacto por su correo tarda & \bl[1.8cm] ms & \bl[1.8cm] ms \\
\bottomrule
\end{tabular}
\end{center}
\vspace{0.15cm}
\small \textbf{Para probarlo}, en cada carpeta: \texttt{npm install}, \texttt{npm run crear} y \texttt{npm run dev}.
\end{frame}
% ---------------------------------------------------------------------
\begin{frame}{\ins{Práctica 1} ¿Índice sí o no?}
La tabla \tbl{alumno} de una universidad tiene \textbf{50\,000 filas}.
\vspace{0.3cm}
\begin{center}
\small
\renewcommand{\arraystretch}{1.15}
\begin{tabular}{>{\raggedright\arraybackslash}p{8.4cm}>{\raggedright\arraybackslash}p{4.0cm}}
\toprule
\bfseries Situación & \bfseries ¿Creo un índice? ¿Por qué? \\
\midrule
Secretaría busca todos los días por \texttt{apellido}. & \bl[3.8cm] \\[10pt]
La columna \texttt{observacion} casi nunca se consulta. & \bl[3.8cm] \\[10pt]
Otra tabla, \tbl{dia\_semana}, tiene solo 7 filas. & \bl[3.8cm] \\
\bottomrule
\end{tabular}
\end{center}
\vspace{0.35cm}
\small \textbf{Para escribir:} la instrucción que crea el índice que sí hace falta.
\end{frame}
% =======================================================================
\seccion[0]{Valores nulos}
% ---------------------------------------------------------------------
\begin{frame}{\texttt{NULL}: el dato que falta}
\begin{columns}[c,onlytextwidth]
\column{.56\textwidth}
\tablaempleado
\column{.40\textwidth}
\small
\begin{itemize}\setlength\itemsep{6pt}
  \item \nulo{} significa \textbf{«no se sabe»} o \textbf{«no hay dato»}.
  \item Bruno y Elena \textbf{no registraron} teléfono.
  \item Diego tiene \texttt{''}: un \textbf{texto vacío}. Parece lo mismo, pero \textbf{no es} \nulo.
\end{itemize}
\end{columns}
\vspace{0.0cm}
\begin{center}
\small
\renewcommand{\arraystretch}{1.05}
\begin{tabular}{>{\ttfamily\bfseries\raggedright\arraybackslash}p{2.2cm}>{\raggedright\arraybackslash}p{9.4cm}}
\toprule
\normalfont\bfseries Valor & \bfseries Qué dice \\
\midrule
0 & un número: «tiene cero» \\
'' & un texto sin letras: «lo dejó en blanco» \\
NULL & \textbf{ningún valor}: «no lo sabemos» \\
\bottomrule
\end{tabular}
\end{center}
\vspace{-0.15cm}
\analogia{En un examen: un \textbf{cero} es una nota. Un casillero \textbf{sin nota} es otra cosa: todavía no se sabe cuánto sacó.}
\end{frame}
% ---------------------------------------------------------------------
\begin{frame}{\texttt{IS NULL}: preguntar si falta el dato}
\small \textbf{Pregunta:} ¿qué empleados \textbf{no tienen} teléfono registrado?
\vspace{0.05cm}
\begin{columns}[c,onlytextwidth]
\column{.56\textwidth}
\begin{sqlcode}[\textwidth]
\scriptsize
\kw{SELECT} nombre, telefono\\
\kw{FROM} empleado\\
\kw{WHERE} telefono \nuevo{\kw{IS NULL}};
\end{sqlcode}
\column{.38\textwidth}
\res{ll}{\cab{nombre} & \cab{telefono}}{%
  Bruno & \nulo \\
  Elena & \nulo \\}
\end{columns}
\vspace{0.2cm}
\begin{itemize}\setlength\itemsep{6pt}
  \item Lo contrario es \texttt{IS NOT NULL}: los que \textbf{sí} tienen un valor.
  \item Diego \textbf{no sale}: su teléfono es \texttt{''}, y eso es un valor.
\end{itemize}
\vspace{0.15cm}
\ojo{\texttt{WHERE telefono = NULL} \textbf{no devuelve nada}. Comparar con algo desconocido nunca da «verdadero». Para los nulos se usa siempre \texttt{IS NULL}.}
\end{frame}
% ---------------------------------------------------------------------
\begin{frame}{\texttt{COALESCE}: un valor de repuesto}
\begin{columns}[c,onlytextwidth]
\column{.58\textwidth}
\begin{sqlcode}[\textwidth]
\scriptsize
\kw{SELECT} nombre,\\
\ind \nuevo{\kw{COALESCE}(telefono, 'sin teléfono')} \kw{AS} contacto\\
\kw{FROM} empleado;
\end{sqlcode}
\column{.38\textwidth}
\res{ll}{\cab{nombre} & \cab{contacto}}{%
  Ana & 987111222 \\
  Bruno & \mal{sin teléfono} \\
  Carla & 955333444 \\
  Diego & '' \\
  Elena & \mal{sin teléfono} \\}
\end{columns}
\vspace{0.25cm}
\begin{itemize}\setlength\itemsep{6pt}
  \item \texttt{COALESCE(a, b)} devuelve \texttt{a}; y si \texttt{a} es \nulo, devuelve \texttt{b}.
  \item Acepta más valores: \texttt{COALESCE(a, b, c)} devuelve \textbf{el primero que no sea} \nulo.
  \item Solo cambia lo que \textbf{se muestra}: la tabla sigue guardando \nulo.
\end{itemize}
\vspace{0.1cm}
\analogia{Es el \textbf{suplente}: si el titular no vino, entra el que sigue en la lista.}
\end{frame}
% ---------------------------------------------------------------------
\begin{frame}{\texttt{NULLIF}: convertir un valor en \texttt{NULL}}
\begin{columns}[c,onlytextwidth]
\column{.58\textwidth}
\begin{sqlcode}[\textwidth]
\scriptsize
\kw{SELECT} nombre,\\
\ind \nuevo{\kw{NULLIF}(telefono, '')} \kw{AS} telefono\\
\kw{FROM} empleado;
\end{sqlcode}
\column{.38\textwidth}
\res{ll}{\cab{nombre} & \cab{telefono}}{%
  Ana & 987111222 \\
  Bruno & \nulo \\
  Carla & 955333444 \\
  Diego & \mal{\nulo} \\
  Elena & \nulo \\}
\end{columns}
\vspace{0.2cm}
\begin{itemize}\setlength\itemsep{5pt}
  \item \texttt{NULLIF(a, b)} devuelve \nulo{} si \texttt{a} \textbf{es igual} a \texttt{b}. Si no, devuelve \texttt{a}.
  \item Juntos: \texttt{COALESCE(NULLIF(telefono, ''), 'sin teléfono')} y Diego también sale «sin teléfono».
\end{itemize}
\vspace{0.1cm}
\enlarealidad{El uso más común es \textbf{evitar dividir entre cero}: \texttt{total / NULLIF(cantidad, 0)} da \nulo{} en vez de un error.}
\end{frame}
% ---------------------------------------------------------------------
\begin{frame}{\ins{Práctica 2} Valores nulos en la tienda}
\begin{columns}[c,onlytextwidth]
\column{.46\textwidth}
\tablaproducto
\column{.50\textwidth}
\small El \texttt{stock} en \nulo{} quiere decir que \textbf{todavía no se contó}.
\vspace{0.15cm}

El \texttt{stock} en \texttt{0} quiere decir que \textbf{se contó y no hay}.
\end{columns}
\vspace{0.25cm}
\begin{center}
\small
\renewcommand{\arraystretch}{1.1}
\begin{tabular}{>{\raggedright\arraybackslash}p{8.8cm}>{\raggedright\arraybackslash}p{3.4cm}}
\toprule
\bfseries Pregunta & \bfseries Respuesta \\
\midrule
¿Cuántas filas devuelve \texttt{WHERE stock IS NULL}? & \bl[3.2cm] \\[8pt]
¿Cuántas filas devuelve \texttt{WHERE stock = NULL}? & \bl[3.2cm] \\[8pt]
¿Qué muestra \texttt{COALESCE(stock, 0)} para el Mouse? & \bl[3.2cm] \\[8pt]
¿Qué muestra \texttt{NULLIF(stock, 0)} para la Mesa? & \bl[3.2cm] \\
\bottomrule
\end{tabular}
\end{center}
\vspace{0.2cm}
\small \textbf{Para escribir:} la consulta que muestra el nombre y el stock, con el texto \texttt{'por contar'} donde falte.
\end{frame}
% =======================================================================
\seccion[0]{Expresiones}
% ---------------------------------------------------------------------
\begin{frame}{Una expresión: calcular dentro del \texttt{SELECT}}
\begin{center}
\begin{tikzpicture}[y=0.82cm]
  \renewcommand{\piezax}{6.0}
  \pieza{2}{\kw{SELECT} nombre,}{una columna \textbf{tal cual} está en la tabla}
  \pieza{1}{\ind \nuevo{sueldo * 12} \kw{AS} sueldo\_anual}{una \textbf{expresión}: un cálculo, con su nombre}
  \pieza{0}{\kw{FROM} empleado;}{la tabla \textbf{no cambia}}
\end{tikzpicture}
\end{center}
\vspace{0.1cm}
\begin{columns}[c,onlytextwidth]
\column{.34\textwidth}
\res{lr}{\cab{nombre} & \cab{sueldo\_anual}}{%
  Ana & 30000 \\
  Bruno & 21600 \\
  Carla & 50400 \\
  Diego & 37200 \\
  Elena & 18000 \\}
\column{.62\textwidth}
\small
\begin{itemize}\setlength\itemsep{6pt}
  \item Una \textbf{expresión} combina columnas, valores y operadores: \texttt{+} \ \texttt{-} \ \texttt{*} \ \texttt{/}
  \item El cálculo se hace \textbf{en cada fila}.
  \item \texttt{AS} le pone nombre a la columna nueva.
  \item La columna \texttt{sueldo\_anual} \textbf{no se guarda}: se calcula al consultar.
\end{itemize}
\end{columns}
\end{frame}
% ---------------------------------------------------------------------
\begin{frame}{Expresiones con fechas: la antigüedad}
\small \textbf{Pregunta:} ¿cuántos años lleva cada empleado en la empresa?
\vspace{0.05cm}
\begin{columns}[c,onlytextwidth]
\column{.64\textwidth}
\begin{sqlcode}[\textwidth]
\scriptsize
\kw{SELECT} nombre,\\
\ind \nuevo{\kw{CAST}(strftime('\%Y', 'now') \kw{AS} INTEGER) - anio\_ingreso}\\
\ind \ind \kw{AS} antiguedad\\
\kw{FROM} empleado;
\end{sqlcode}
\column{.32\textwidth}
\res{lr}{\cab{nombre} & \cab{antiguedad}}{%
  Ana & 8 \\
  Bruno & 5 \\
  Carla & 11 \\
  Diego & 3 \\
  Elena & 6 \\}
\end{columns}
\vspace{0.15cm}
\begin{center}
\small
\renewcommand{\arraystretch}{1.25}
\begin{tabular}{>{\ttfamily\raggedright\arraybackslash}p{5.2cm}>{\raggedright\arraybackslash}p{7.2cm}}
\toprule
\normalfont\bfseries Pieza & \bfseries Qué hace \\
\midrule
strftime('\%Y', 'now') & el año de hoy, como \textbf{texto}: \texttt{'2026'} \\
CAST(... AS INTEGER) & lo convierte en \textbf{número}: \texttt{2026} \\
- anio\_ingreso & le resta el año en que entró cada uno \\
\bottomrule
\end{tabular}
\end{center}
\vspace{0.05cm}
\small Así se escribe en \textbf{SQLite}. En MySQL la misma idea es \texttt{YEAR(CURDATE()) - anio\_ingreso}.
\end{frame}
% ---------------------------------------------------------------------
\begin{frame}{Concatenar: pegar textos con \texttt{||}}
\begin{columns}[c,onlytextwidth]
\column{.60\textwidth}
\begin{sqlcode}[\textwidth]
\scriptsize
\kw{SELECT} \nuevo{nombre || ' (' || area || ')'} \kw{AS} info\\
\kw{FROM} empleado;
\end{sqlcode}
\column{.34\textwidth}
\res{l}{\cab{info}}{%
  Ana (Ventas) \\
  Bruno (Ventas) \\
  Carla (Sistemas) \\
  Diego (Sistemas) \\
  Elena (Almacén) \\}
\end{columns}
\vspace{0.25cm}
\begin{itemize}\setlength\itemsep{6pt}
  \item \texttt{||} \textbf{pega} un texto a continuación del otro.
  \item Lo que va entre comillas se escribe \textbf{tal cual}: aquí, los paréntesis y el espacio.
  \item También pega números: \texttt{nombre || ' - ' || sueldo} da \texttt{Ana - 2500}.
\end{itemize}
\vspace{0.1cm}
\ojo{Si \textbf{una} de las piezas es \nulo, \textbf{todo} el resultado es \nulo. Se arregla con \texttt{COALESCE(telefono, 'sin teléfono')}.}
\vspace{0.05cm}
\small En MySQL se usa \texttt{CONCAT(nombre, ' (', area, ')')}.
\end{frame}
% ---------------------------------------------------------------------
\begin{frame}{\texttt{CASE WHEN}: decidir fila por fila}
\begin{columns}[c,onlytextwidth]
\column{.56\textwidth}
\begin{sqlcode}[\textwidth]
\scriptsize
\kw{SELECT} nombre, sueldo,\\
\ind \nuevo{\kw{CASE}}\\
\ind \ind \nuevo{\kw{WHEN} sueldo >= 3000 \kw{THEN} 'alto'}\\
\ind \ind \nuevo{\kw{WHEN} sueldo >= 2000 \kw{THEN} 'medio'}\\
\ind \ind \nuevo{\kw{ELSE} 'bajo'}\\
\ind \nuevo{\kw{END}} \kw{AS} nivel\\
\kw{FROM} empleado;
\end{sqlcode}
\column{.40\textwidth}
\res{lrl}{\cab{nombre} & \cab{sueldo} & \cab{nivel}}{%
  Ana & 2500 & medio \\
  Bruno & 1800 & bajo \\
  Carla & 4200 & alto \\
  Diego & 3100 & alto \\
  Elena & 1500 & bajo \\}
\end{columns}
\vspace{0.2cm}
\begin{itemize}\setlength\itemsep{5pt}
  \item Las condiciones se revisan \textbf{en orden}: se queda con \textbf{la primera} que se cumple.
  \item \texttt{ELSE} atrapa a todos los demás. Sin \texttt{ELSE}, a esos les toca \nulo.
  \item Siempre se cierra con \texttt{END}.
\end{itemize}
\vspace{0.05cm}
\analogia{Es la regla del profesor: «de 18 para arriba, sobresaliente; de 11 para arriba, aprobado; \textbf{si no}, desaprobado».}
\end{frame}
% ---------------------------------------------------------------------
\begin{frame}{Funciones de texto}
\small Una \textbf{función} recibe un valor y devuelve otro. Estas trabajan con \textbf{texto}.
\begin{center}
\small
\renewcommand{\arraystretch}{1.3}
\begin{tabular}{>{\ttfamily\bfseries\raggedright\arraybackslash}p{1.7cm}>{\raggedright\arraybackslash}p{4.3cm}>{\ttfamily\raggedright\arraybackslash}p{4.7cm}>{\ttfamily\raggedright\arraybackslash}p{1.5cm}}
\toprule
\normalfont\bfseries Función & \bfseries Qué hace & \normalfont\bfseries Con \texttt{nombre = 'Carla'} & \normalfont\bfseries Resultado \\
\midrule
UPPER & pasa todo a mayúsculas & UPPER(nombre) & CARLA \\
LOWER & pasa todo a minúsculas & LOWER(nombre) & carla \\
LENGTH & cuenta los caracteres & LENGTH(nombre) & 5 \\
SUBSTR & saca un pedazo: desde dónde y cuántos & SUBSTR(nombre, 1, 3) & Car \\
REPLACE & cambia un texto por otro & REPLACE(nombre, 'a', 'o') & Corlo \\
\bottomrule
\end{tabular}
\end{center}
\vspace{-0.1cm}
\begin{sqlcode}[.80\textwidth]
\scriptsize
\kw{SELECT} \kw{UPPER}(nombre) \kw{AS} nombre, \kw{SUBSTR}(area, 1, 3) \kw{AS} sigla \ \ \com{CARLA | Sis}\\
\kw{FROM} empleado;
\end{sqlcode}
\vspace{-0.1cm}
\small En \texttt{SUBSTR} las posiciones empiezan en \textbf{1}, no en 0.
\end{frame}
% ---------------------------------------------------------------------
\begin{frame}{Funciones numéricas}
\begin{columns}[c,onlytextwidth]
\column{.60\textwidth}
\begin{sqlcode}[\textwidth]
\scriptsize
\kw{SELECT} nombre,\\
\ind \nuevo{\kw{ROUND}(sueldo / 30.0, 2)} \kw{AS} sueldo\_diario,\\
\ind \nuevo{\kw{ABS}(sueldo - 2620)} \kw{AS} distancia\\
\kw{FROM} empleado;
\end{sqlcode}
\column{.38\textwidth}
\res{lrr}{\cab{nombre} & \cab{sueldo\_diario} & \cab{distancia}}{%
  Ana & 83.33 & 120 \\
  Bruno & 60.0 & 820 \\
  Carla & 140.0 & 1580 \\
  Diego & 103.33 & 480 \\
  Elena & 50.0 & 1120 \\}
\end{columns}
\vspace{0.15cm}
\begin{center}
\small
\renewcommand{\arraystretch}{1.25}
\begin{tabular}{>{\ttfamily\bfseries\raggedright\arraybackslash}p{1.5cm}>{\raggedright\arraybackslash}p{5.6cm}>{\ttfamily\raggedright\arraybackslash}p{3.4cm}>{\ttfamily\raggedright\arraybackslash}p{1.5cm}}
\toprule
\normalfont\bfseries Función & \bfseries Qué hace & \normalfont\bfseries Ejemplo & \normalfont\bfseries Resultado \\
\midrule
ROUND & redondea a los decimales que pidas & ROUND(83.3333, 2) & 83.33 \\
ABS & quita el signo: el valor absoluto & ABS(-820) & 820 \\
\bottomrule
\end{tabular}
\end{center}
\vspace{0.05cm}
\small \texttt{2620} es el sueldo promedio: \texttt{distancia} dice \textbf{cuánto se aleja} cada uno, por arriba o por abajo.
\vspace{0.1cm}
\ojo{Entre dos enteros, \texttt{/} \textbf{bota los decimales}: \texttt{2500 / 30} da \texttt{83}. Por eso se divide entre \texttt{30.0}.}
\end{frame}
% ---------------------------------------------------------------------
\begin{frame}{\ins{Práctica 3} Expresiones en la tienda}
\begin{columns}[c,onlytextwidth]
\column{.44\textwidth}
\tablaproducto
\column{.54\textwidth}
\begin{sqlcode}[\textwidth]
\scriptsize
\kw{SELECT} \kw{UPPER}(nombre) \kw{AS} nombre,\\
\ind \kw{ROUND}(precio * 1.18, 1) \kw{AS} precio\_final,\\
\ind \kw{CASE}\\
\ind \ind \kw{WHEN} precio >= 500 \kw{THEN} 'caro'\\
\ind \ind \kw{ELSE} 'económico'\\
\ind \kw{END} \kw{AS} tipo\\
\kw{FROM} producto;
\end{sqlcode}
\end{columns}
\vspace{0.1cm}
\begin{center}
\small
\renewcommand{\arraystretch}{1.1}
\begin{tabular}{>{\raggedright\arraybackslash}p{8.8cm}>{\raggedright\arraybackslash}p{3.4cm}}
\toprule
\bfseries Pregunta & \bfseries Respuesta \\
\midrule
¿Qué sale en \texttt{nombre} para la fila 3? & \bl[3.2cm] \\[5pt]
¿Cuánto es el \texttt{precio\_final} del Mouse? & \bl[3.2cm] \\[5pt]
¿Cuántos productos salen como \texttt{'caro'}? & \bl[3.2cm] \\
\bottomrule
\end{tabular}
\end{center}
\vspace{0.1cm}
\small \textbf{Para escribir:} la consulta que muestra \texttt{Laptop [Cómputo]}, \texttt{Mouse [Cómputo]}\ldots{} en una columna llamada \texttt{etiqueta}.
\end{frame}
% =======================================================================
\seccion[0]{Agregaciones}
% ---------------------------------------------------------------------
\begin{frame}{Agregaciones: muchas filas, un solo valor}
\begin{columns}[c,onlytextwidth]
\column{.54\textwidth}
\tablaempleado
\column{.44\textwidth}
\small
\begin{itemize}\setlength\itemsep{6pt}
  \item Una expresión calcula \textbf{fila por fila}.
  \item Una \textbf{agregación} toma \textbf{toda una columna} y devuelve \textbf{un solo valor}.
\end{itemize}
\end{columns}
\vspace{0.0cm}
\begin{center}
\small
\renewcommand{\arraystretch}{1.08}
\begin{tabular}{>{\ttfamily\bfseries\raggedright\arraybackslash}p{1.5cm}>{\raggedright\arraybackslash}p{4.6cm}>{\ttfamily\raggedright\arraybackslash}p{3.4cm}>{\ttfamily\raggedleft\arraybackslash}p{1.7cm}}
\toprule
\normalfont\bfseries Función & \bfseries Qué responde & \normalfont\bfseries Ejemplo & \normalfont\bfseries Resultado \\
\midrule
COUNT & ¿cuántas filas hay? & COUNT(*) & 5 \\
SUM & ¿cuánto suman? & SUM(sueldo) & 13100 \\
AVG & ¿cuál es el promedio? & AVG(sueldo) & 2620.0 \\
MIN & ¿cuál es el menor? & MIN(sueldo) & 1500 \\
MAX & ¿cuál es el mayor? & MAX(sueldo) & 4200 \\
\bottomrule
\end{tabular}
\end{center}
\vspace{-0.1cm}
\begin{sqlcode}[.80\textwidth]
\scriptsize
\kw{SELECT} \kw{COUNT}(*) \kw{AS} cantidad, \kw{AVG}(sueldo) \kw{AS} sueldo\_promedio \kw{FROM} empleado;
\end{sqlcode}
\end{frame}
% ---------------------------------------------------------------------
\begin{frame}{Las agregaciones y los \texttt{NULL}}
\begin{columns}[c,onlytextwidth]
\column{.56\textwidth}
\begin{sqlcode}[\textwidth]
\scriptsize
\kw{SELECT} \kw{COUNT}(*) \kw{AS} empleados,\\
\ind \nuevo{\kw{COUNT}(telefono)} \kw{AS} con\_telefono\\
\kw{FROM} empleado;
\end{sqlcode}
\column{.40\textwidth}
\res{rr}{\cab{empleados} & \cab{con\_telefono}}{%
  5 & \mal{3} \\}
\end{columns}
\vspace{0.3cm}
\begin{center}
\small
\renewcommand{\arraystretch}{1.3}
\begin{tabular}{>{\ttfamily\bfseries\raggedright\arraybackslash}p{3.4cm}>{\raggedright\arraybackslash}p{8.6cm}}
\toprule
\normalfont\bfseries Se escribe & \bfseries Cuenta \\
\midrule
COUNT(*) & \textbf{todas} las filas \\
COUNT(telefono) & solo las filas donde \texttt{telefono} \textbf{no es} \nulo \\
\bottomrule
\end{tabular}
\end{center}
\vspace{0.15cm}
\begin{itemize}\setlength\itemsep{6pt}
  \item Bruno y Elena no se cuentan. Diego \textbf{sí}: su \texttt{''} es un valor.
  \item \texttt{SUM}, \texttt{AVG}, \texttt{MIN} y \texttt{MAX} también \textbf{saltan} los \nulo: no los toman como cero.
\end{itemize}
\end{frame}
% ---------------------------------------------------------------------
\begin{frame}{\texttt{GROUP BY}: un resultado por grupo}
\small \textbf{Pregunta:} ¿cuántos empleados hay y cuál es el sueldo promedio \textbf{en cada área}?
\vspace{0.05cm}
\begin{columns}[c,onlytextwidth]
\column{.50\textwidth}
\begin{sqlcode}[\textwidth]
\scriptsize
\kw{SELECT} area,\\
\ind \kw{COUNT}(*) \kw{AS} cantidad,\\
\ind \kw{AVG}(sueldo) \kw{AS} sueldo\_promedio\\
\kw{FROM} empleado\\
\nuevo{\kw{GROUP BY} area};
\end{sqlcode}
\column{.48\textwidth}
\res{lrr}{\cab{area} & \cab{cantidad} & \cab{sueldo\_promedio}}{%
  Almacén & 1 & 1500.0 \\
  Sistemas & 2 & 3650.0 \\
  Ventas & 2 & 2150.0 \\}
\end{columns}
\vspace{0.2cm}
\begin{enumerate}\setlength\itemsep{5pt}
  \item Se \textbf{juntan} las filas que tienen la misma \texttt{area}: salen tres grupos.
  \item La agregación se calcula \textbf{dentro de cada grupo}.
  \item El resultado tiene \textbf{una fila por grupo}.
\end{enumerate}
\vspace{0.1cm}
\ojo{En el \texttt{SELECT} solo van las columnas del \texttt{GROUP BY} y las agregaciones. Pedir también \texttt{nombre} no tiene sentido: ¿el de cuál de los dos de Ventas?}
\end{frame}
% ---------------------------------------------------------------------
\begin{frame}{\texttt{HAVING}: filtrar los grupos}
\small \textbf{Pregunta:} ¿qué áreas tienen \textbf{dos o más} empleados?
\vspace{0.05cm}
\begin{columns}[c,onlytextwidth]
\column{.56\textwidth}
\begin{sqlcode}[\textwidth]
\scriptsize
\kw{SELECT} area, \kw{COUNT}(*) \kw{AS} cantidad\\
\kw{FROM} empleado\\
\kw{GROUP BY} area\\
\nuevo{\kw{HAVING} \kw{COUNT}(*) >= 2};
\end{sqlcode}
\column{.38\textwidth}
\res{lr}{\cab{area} & \cab{cantidad}}{%
  Sistemas & 2 \\
  Ventas & 2 \\}
\end{columns}
\vspace{-0.05cm}
\begin{center}
\small
\renewcommand{\arraystretch}{1.1}
\begin{tabular}{>{\ttfamily\bfseries\raggedright\arraybackslash}p{1.7cm}>{\raggedright\arraybackslash}p{4.4cm}>{\raggedright\arraybackslash}p{5.8cm}}
\toprule
& \bfseries Qué filtra & \bfseries Cuándo \\
\midrule
WHERE & \textbf{filas} sueltas & \textbf{antes} de agrupar \\
HAVING & \textbf{grupos} completos & \textbf{después} de agrupar \\
\bottomrule
\end{tabular}
\end{center}
\vspace{-0.15cm}
\ojo{\texttt{WHERE COUNT(*) >= 2} da error: cuando actúa el \texttt{WHERE}, los grupos \textbf{todavía no existen}.}
\end{frame}
% ---------------------------------------------------------------------
\begin{frame}{Todo junto: el orden de una consulta}
\begin{center}
\begin{tikzpicture}[y=0.74cm]
  \renewcommand{\piezax}{7.6}
  \pieza{5}{\kw{SELECT} area, \kw{AVG}(sueldo) \kw{AS} sueldo\_promedio}{qué se muestra de cada grupo}
  \pieza{4}{\kw{FROM} empleado}{de qué tabla}
  \pieza{3}{\kw{WHERE} anio\_ingreso >= 2018}{qué \textbf{filas} entran}
  \pieza{2}{\kw{GROUP BY} area}{cómo se \textbf{agrupan}}
  \pieza{1}{\kw{HAVING} \kw{AVG}(sueldo) > 2000}{qué \textbf{grupos} se quedan}
  \pieza{0}{\kw{ORDER BY} sueldo\_promedio \kw{DESC};}{en qué orden salen}
\end{tikzpicture}
\end{center}
\vspace{0.1cm}
\begin{columns}[c,onlytextwidth]
\column{.36\textwidth}
\res{lr}{\cab{area} & \cab{sueldo\_promedio}}{%
  Sistemas & 3100.0 \\
  Ventas & 2150.0 \\}
\column{.60\textwidth}
\small
\begin{itemize}\setlength\itemsep{5pt}
  \item El \texttt{WHERE} deja fuera a \textbf{Carla}, que entró en 2015.
  \item Quedan tres grupos: Ventas, Sistemas y Almacén.
  \item El \texttt{HAVING} deja fuera a \textbf{Almacén}: su promedio es 1500.
\end{itemize}
\end{columns}
\vspace{0.15cm}
\small Las cláusulas se escriben \textbf{siempre en este orden}.
\end{frame}
% ---------------------------------------------------------------------
\begin{frame}{\ins{Práctica 4} Haz tú el \texttt{GROUP BY}, a mano}
\begin{columns}[c,onlytextwidth]
\column{.44\textwidth}
\tablaproducto
\column{.54\textwidth}
\begin{sqlcode}[\textwidth]
\scriptsize
\kw{SELECT} categoria,\\
\ind \kw{COUNT}(*) \kw{AS} cantidad,\\
\ind \kw{MAX}(precio) \kw{AS} precio\_mayor\\
\kw{FROM} producto\\
\kw{GROUP BY} categoria;
\end{sqlcode}
\end{columns}
\vspace{0.15cm}
\begin{columns}[c,onlytextwidth]
\column{.50\textwidth}
\small
\begin{enumerate}\setlength\itemsep{7pt}
  \item \textbf{Agrupa:} en la tabla, encierra juntas las filas de la misma \texttt{categoria}.
  \item \textbf{Completa} el resultado: una fila por cada grupo.
  \item \textbf{Tacha} la fila que desaparece si agregamos \texttt{HAVING COUNT(*) >= 2}.
\end{enumerate}
\column{.46\textwidth}
\centering{\scriptsize\color{utpGris}resultado}\\[3pt]
\footnotesize
\renewcommand{\arraystretch}{1.5}
\begin{tabular}{lll}
\toprule
\ttfamily\cab{categoria} & \ttfamily\cab{cantidad} & \ttfamily\cab{precio\_mayor} \\
\midrule
\bl[1.7cm] & \bl[1.2cm] & \bl[1.5cm] \\
\bl[1.7cm] & \bl[1.2cm] & \bl[1.5cm] \\
\bl[1.7cm] & \bl[1.2cm] & \bl[1.5cm] \\
\bottomrule
\end{tabular}
\end{columns}
\vspace{0.4cm}
\small Trabaja \textbf{en pareja}: uno agrupa, el otro completa. Luego comparen con la pareja de al lado.
\end{frame}
% ---------------------------------------------------------------------
\begin{frame}[plain]
  \begin{tikzpicture}[remember picture,overlay]
    \node[anchor=north west,font=\bfseries\fontsize{24}{28}\selectfont]
      at ([xshift=12mm,yshift=-9mm]current page.north west) {Preguntas};
    \node[anchor=north east,inner sep=0pt]
      at ([xshift=-10mm,yshift=-8mm]current page.north east)
      {\includegraphics[width=43mm]{utp-logo.png}};
    \node[anchor=center,inner sep=0pt]
      at ([yshift=-5mm]current page.center)
      {\includegraphics[width=78mm]{preguntas-ilustracion.png}};
  \end{tikzpicture}
\end{frame}
% ---------------------------------------------------------------------
\begin{frame}{Cierre: lo de hoy en una tabla}
\begin{center}
\renewcommand{\arraystretch}{1.4}
\begin{tabular}{>{\raggedright\bfseries}p{4.2cm}>{\raggedright\arraybackslash}p{8.6cm}}
\toprule
& \bfseries En una frase \\
\midrule
Índice & Una lista ordenada que dice dónde está cada fila. \\
\texttt{IS NULL} & Pregunta si falta el dato. Nunca \texttt{= NULL}. \\
\texttt{COALESCE} / \texttt{NULLIF} & Poner un repuesto al \nulo; convertir un valor en \nulo. \\
Expresión & Un cálculo en el \texttt{SELECT}, fila por fila, con su \texttt{AS}. \\
\texttt{CASE WHEN} & Decide un valor según una condición. \\
Agregación & Muchas filas, un solo valor: \texttt{COUNT}, \texttt{SUM}, \texttt{AVG}, \texttt{MIN}, \texttt{MAX}. \\
\texttt{GROUP BY} / \texttt{HAVING} & Un resultado por grupo; quedarse solo con algunos grupos. \\
\bottomrule
\end{tabular}
\end{center}
\vspace{0.2cm}
\small Para no confundirlos: \texttt{WHERE} filtra \textbf{filas}, \texttt{HAVING} filtra \textbf{grupos}.
\end{frame}
% ---------------------------------------------------------------------
\begin{frame}{Referencias}
\footnotesize
Coronel, C., Morris, S., \& Rob, P. (2011). \textit{Bases de datos: diseño, implementación y administración} (9.ª ed.). Cengage Learning.

\vspace{0.2cm}
Elmasri, R., \& Navathe, S. (2007). \textit{Fundamentos de sistemas de bases de datos} (5.ª ed.). Addison Wesley.

\vspace{0.2cm}
Silberschatz, A., Korth, H. F., \& Sudarshan, S. (2020). \textit{Database system concepts} (7th ed.). McGraw-Hill.

\vspace{0.2cm}
SQLite. (s.\,f.). \textit{Query planning}. \texttt{https://www.sqlite.org/queryplanner.html}

\vspace{0.2cm}
SQLite. (s.\,f.). \textit{Built-in scalar SQL functions}. \texttt{https://www.sqlite.org/lang\_corefunc.html}

\vspace{0.2cm}
SQLite. (s.\,f.). \textit{Built-in aggregate functions}. \texttt{https://www.sqlite.org/lang\_aggfunc.html}
\end{frame}
% ---------------------------------------------------------------------
% Cierre UTP
{
  \setbeamertemplate{background canvas}{%
    \includegraphics[width=\paperwidth,height=\paperheight]{gracias-final.pdf}%
  }
  \setbeamertemplate{footline}{}
  \begin{frame}[plain]
  \end{frame}
}

\end{document}
